\section{Results}
\label{sec:results}

\subsection{Count-Degradation Curves}
\label{sec:count_curves}

We first examine how extraction performance changes as the number of types in the prompt schema increases (Experiment~1). Figure~\ref{fig:count_curves} plots per-task metrics across count levels.

\paragraph{NER (Few-NERD).}
Recall declines monotonically from 0.72 at count${}\,=\,1$ to 0.42 at count${}\,=\,64$, a 42\% relative drop. Precision moves in the opposite direction, rising from 0.11 to 0.52, reflecting a shift from over-prediction to under-prediction as more candidate types compete in the prompt. F1 exhibits an inverted-U shape, peaking at intermediate counts where the precision--recall trade-off is most balanced.

\paragraph{RE (FewRel).}
Accuracy drops monotonically from 85.7\% at count${}\,=\,1$ to 50.4\% at count${}\,=\,80$, a near-halving when the model must discriminate among all 80 relation types simultaneously.

\paragraph{ED (MAVEN).}
In contrast to NER and RE, MAVEN relaxed F1 \emph{increases} from 0.106 at count${}\,=\,1$ to 0.247 at count${}\,=\,44$. This reversal arises from three properties of the MAVEN setup: (i)~trigger-level relaxed matching, which requires only token overlap rather than exact span match; (ii)~MAVEN's 168 event types derive from the FrameNet ontology, forming a same-domain hierarchy where co-presented types provide disambiguating context; and (iii)~extreme class imbalance---rare event types benefit from being presented alongside related types that prime the model toward the correct semantic region. We return to this mechanism in \S\ref{sec:centrality_confound}.

\paragraph{Summary.}
Schema size degrades extraction quality for NER and RE (Channel~1: count $\to$ recall decline), but this pattern does not hold universally. MAVEN's reversed trend demonstrates that task structure and evaluation protocol mediate the count--performance relationship.

\subsection{Pairwise Confusability and Type Confusion}
\label{sec:pairwise}

We test whether pairwise SCS predicts type-level confusion rates. For each type pair $(t_i, t_j)$, we compute the confusion rate---the fraction of $t_i$ gold instances misclassified as $t_j$---and correlate it with $\text{SCS}(t_i, t_j)$. Table~\ref{tab:pairwise} reports Pearson correlations with bootstrap 95\% CIs.

\begin{table}[t]
\centering
\small
\begin{tabular}{llrrrrr}
\toprule
\textbf{Dataset} & \textbf{Embedding} & \textbf{$n$ pairs} & \textbf{$r$} & \textbf{$r^2$} & \textbf{95\% CI} & \textbf{$p$} \\
\midrule
Few-NERD & Qwen3 & 1{,}953 & 0.112 & 0.013 & {[}0.076, 0.156{]} & ${<}10^{-12}$ \\
MAVEN & MiniLM & 11{,}026 & 0.129 & 0.017 & {[}0.104, 0.156{]} & ${<}10^{-10}$ \\
\bottomrule
\end{tabular}
\caption{Pairwise SCS--confusion correlation. $r^2$ indicates the fraction of variance in type confusion explained by embedding similarity. Both correlations are modest in magnitude but robust across datasets and embedding models ($p < 10^{-10}$).}
\label{tab:pairwise}
\end{table}

The correlations are modest: $r^2 \approx 1.3\text{--}1.7\%$, meaning pairwise confusability explains roughly 1.5\% of the variance in type confusion. Other factors---entity frequency, prompt position effects \cite{lost-in-middle-2024}, and type definition quality---jointly account for the remainder. Nevertheless, the signal is robust: both correlations survive bootstrap resampling with CIs bounded well above zero, and replicate across two unrelated embedding spaces (Qwen3 hidden states and MiniLM sentence embeddings).

Spearman rank correlations from the SCS validation experiment (exp-015) corroborate these findings. Across all type pairs, Qwen3 yields $\rho = 0.168$ ($p = 7.3 \times 10^{-14}$) and MiniLM yields $\rho = 0.144$ ($p = 1.7 \times 10^{-10}$). Restricting to pairs with nonzero confusion strengthens the association: $\rho = 0.272$ ($p = 5 \times 10^{-5}$) for Qwen3 and $\rho = 0.306$ ($p = 4.6 \times 10^{-6}$) for MiniLM, indicating that among types the model actually confuses, embedding similarity is a stronger predictor.

We frame these results as identifying one factor among several that drive type confusion, but the first to be empirically isolated in a controlled zero-shot IE setting.

\subsection{BCS-Controlled Experiments}
\label{sec:bcs_controlled}

To move from correlational to interventional evidence, we fix the type count and vary BCS. Table~\ref{tab:bcs_controlled} reports FewRel accuracy under low-BCS and high-BCS batches at two schema sizes; Figure~\ref{fig:bcs_scaling} visualizes the scaling relationship.

\begin{table}[t]
\centering
\small
\begin{tabular}{lcrrrr}
\toprule
\textbf{Schema} & \textbf{Count} & \textbf{Low-BCS Acc} & \textbf{High-BCS Acc} & \textbf{$\Delta$ (pp)} & \textbf{95\% CI} \\
\midrule
16-rel & 8 & 88.9\% & 79.8\% & $-$9.1 & --- \\
80-rel & 8 & 82.3\% & 51.0\% & $-$31.2 & {[}17.7, 42.7{]} \\
80-rel & 16 & 82.8\% & 51.6\% & $-$31.2 & {[}22.9, 39.6{]} \\
\bottomrule
\end{tabular}
\caption{BCS-controlled FewRel accuracy. Low-BCS batches contain maximally dissimilar types; high-BCS batches contain maximally similar types. The confusability gap widens from 9.1\,pp at 16 relations to 31.2\,pp at 80 relations.}
\label{tab:bcs_controlled}
\end{table}

At the 16-relation scale, switching from low-BCS to high-BCS batches reduces accuracy by 9.1 percentage points. At 80 relations, the gap reaches 31.2\,pp (95\% CI: [17.7, 42.7]). The effect is stable across count levels: at count${}\,=\,16$, the gap remains 31.2\,pp (CI: [22.9, 39.6]).

The larger absolute gap at 80 relations reflects the wider BCS range available in that pool (0.46--0.68 vs.\ 0.65--0.90 for 16-rel), not a stronger per-unit sensitivity. Within-pool analysis confirms that the BCS--accuracy slope is consistent across pools (permutation $p = 0.60$).

\subsection{Multi-Channel Degradation}
\label{sec:multi_channel}

We now decompose extraction errors into an error taxonomy to distinguish the mechanisms behind performance loss. Figure~\ref{fig:error_taxonomy} illustrates the taxonomy; Table~\ref{tab:error_taxonomy} quantifies how each error type scales with BCS.

\begin{table}[t]
\centering
\small
\begin{tabular}{llrrrrrr}
\toprule
 & & \multicolumn{4}{c}{\textbf{BCS Bin}} & & \\
\cmidrule(lr){3-6}
\textbf{Task} & \textbf{Error Type} & \textbf{Solo} & \textbf{Low} & \textbf{Mid} & \textbf{High} & \textbf{$\rho$} & \textbf{$p$} \\
\midrule
NER & IntraSub & 0\% & 1.6\% & 10.6\% & 49.4\% & 1.000 & ${<}0.001$ \\
RE  & IntraSub & 0\% & 4.2\% & 13.0\% & 32.3\% & 0.954 & ${<}10^{-6}$ \\
\midrule
NER & Missed & --- & 12\% & 14\% & 16\% & --- & --- \\
RE  & Missed & --- & 10\% & 12\% & 14\% & --- & --- \\
\bottomrule
\end{tabular}
\caption{Intra-batch substitution (IntraSub) rate by BCS bin, with Spearman $\rho$ against BCS. Missed-entity rates remain flat across bins, indicating that BCS drives type confusion rather than detection failure. Both NER and RE exhibit near-perfect monotonic scaling ($\rho \geq 0.954$).}
\label{tab:error_taxonomy}
\end{table}

Three findings emerge. First, intra-batch substitution scales monotonically with BCS in both tasks: NER IntraSub rises from 0\% (solo) to 49.4\% (high BCS) with $\rho = 1.0$ ($p < 0.001$), and RE IntraSub rises from 0\% to 32.3\% with $\rho = 0.954$ ($p < 10^{-6}$). Second, 100\% of type confusion errors in RE are intra-batch---the model exclusively confuses types present in the current prompt, never hallucinating types outside the schema. Third, the missed-entity rate remains flat at 10--16\% regardless of BCS bin, confirming that confusability drives type confusion (Channel~2) independently of detection failure.

The extreme-batch analysis sharpens this picture. At count${}\,=\,6$, high-BCS batches produce $27.5\times$ more confusion than low-BCS batches under extreme selection, versus $1.7\times$ under random batching. At count${}\,=\,12$, the ratios are $9.9\times$ (extreme) versus $2.4\times$ (random). These multiplicative gaps confirm that BCS captures a causally relevant dimension of schema difficulty, not merely a correlate.

Taken together, these results establish two orthogonal degradation channels: \textbf{Channel~1} (count $\to$ recall decline, \S\ref{sec:count_curves}) and \textbf{Channel~2} (confusability $\to$ type confusion, this section). The channels are mechanistically distinct---Channel~1 reflects capacity limits as more types compete for attention, while Channel~2 reflects semantic similarity making types indistinguishable regardless of count.

\subsection{Centrality Confound on Homogeneous Schemas}
\label{sec:centrality_confound}

Per-type analysis on MAVEN reveals an anomaly: the correlation between mean BCS and per-type error rate is \emph{positive}---Qwen3 $\rho = +0.184$ ($p = 0.025$) and MiniLM $\rho = +0.208$ ($p = 0.011$)---opposite to the negative direction predicted by the confusability hypothesis.

This reversal traces to a property of same-domain schemas documented by \citet{cai-isotropy-clusters-2021}: when all types derive from a single ontology (here, FrameNet), their embeddings collapse into a single cluster in representation space. In this regime, a type's mean BCS conflates two distinct properties: (i)~confusability with neighbors, which should increase errors, and (ii)~centrality within the cluster, which correlates with being a well-represented, high-frequency type that LLMs extract reliably. The centrality signal dominates, flipping the aggregate correlation.

Critically, the pairwise SCS--confusion correlation remains intact: MAVEN MiniLM pairwise $r = 0.129$ (CI: [0.104, 0.156], $p < 10^{-10}$), consistent with the values in Table~\ref{tab:pairwise}. The problem lies in the aggregation step (mean BCS per type), not in the pairwise signal itself.

Few-NERD, whose 66 types span 8 coarse categories across diverse domains, does not exhibit this confound. Its per-type correlation is $\rho = -0.293$ ($p = 0.02$, $n = 63$), in the expected negative direction: types with higher mean confusability have higher error rates. The cross-domain schema spans multiple embedding clusters, breaking the centrality--confusability conflation.

This analysis carries a practical implication: mean BCS is a reliable predictor of schema difficulty only for cross-domain schemas. For same-domain schemas where types form a single semantic cluster, pairwise SCS should be used directly, or the centrality component must be factored out before aggregation.
